\documentclass[11pt,a4paper]{article}

\usepackage[utf8]{inputenc}
\usepackage[T1]{fontenc}
\usepackage{lmodern}
\usepackage[english]{babel}
\usepackage{geometry}
\usepackage{graphicx}
\usepackage{float}
\usepackage{caption}
\usepackage{amsmath}
\usepackage{booktabs}
\usepackage{hyperref}
\usepackage{subcaption}
\usepackage{xcolor}

\geometry{margin=1in}
\raggedbottom
\emergencystretch=3em

% Figures live in the repo's runs/ directory; this file sits in report/6/,
% so resolve \includegraphics paths relative to the repository root.
\graphicspath{{../../}}

% Graceful fallback for figures produced by runs that may not be pulled yet.
\newcommand{\figorbox}[2]{%
  \IfFileExists{../../#1}{\includegraphics[width=#2\textwidth]{#1}}{%
    \fbox{\parbox[c][3.5cm][c]{#2\textwidth}{\centering\itshape
    \textcolor{gray}{[figure not yet generated --- run the eval and pull it]}}}}}

\title{Empirical Analysis of Transformers on Graph Connectivity \\
       \large Part VI: Learning to Reason --- Which Data Teaches a Connectivity Path?}
\author{Edoardo Paccagnella}
\date{\today}

\begin{document}

\maketitle

% =====================================================================
% DRAFT / SKELETON (2026-06-26). No data collected yet. This file lays out
% the framing and the planned experiments for Report VI. Each results
% subsection is a STUB: the % comments record (a) the question, (b) the
% prediction / which thesis it supports, (c) the clean training distribution
% to use, (d) the caption fields required by rule 21. Replace the italic
% "[To run --- data pending]" lines with real tables/figures once the runs
% are in. ONE experiment per chat.
%
% Design rules carried over from istruzioni.md (esp. errors 21, 55-60):
%  - every caption states MODEL + TRAINING SET + TEST SET + METRIC, no analysis;
%  - no code/file/path/flag names in the rendered text;
%  - no 'advisor' attributions, no chronological "we re-ran" story;
%  - NEW for Report VI: NO opaque random "mixed" training stream. Every
%    experiment trains on a single, explicitly named and sensible
%    distribution (ER, cliques, bridged cliques, multipath, ...), or a
%    clearly stated combination. See istruzioni.md sec 13.
% =====================================================================

\begin{abstract}
\noindent \emph{(Provisional.)} The first five reports asked \emph{what} algorithm the standard
connectivity transformer ends up computing (a parallel, distance-bounded matrix powering, not a
sequential traversal; Report~V). This report turns the question around and asks \emph{how it learns
to reason}: which training data let an $L=2$ model acquire a connectivity ``path'' it was never
shown, and where instead the limit is simply not having seen the structure. We organise the study
around two ideas the architecture makes concrete --- \textbf{parallelism} (several routes between
two endpoints) and \textbf{repeated/iterated structure} (a chain of locally-decided links) --- and
around one methodological commitment: every experiment trains on a single, clearly named
distribution, never an opaque random mixture. This is the concluding report of the series, so we
keep the experiment count small and each one pointed at a single claim.
\end{abstract}

\section{Introduction}
\label{sec:intro}

\subsection{What Reports I--V established}
\label{sec:recap}

We treat the following as established (each multi-seed, documented in Reports~I--V); the new study
must stay consistent with them.

\begin{itemize}\setlength\itemsep{0.2em}
  \item \textbf{The data lever does not help the standard transformer generalise}, and whether a
        run finds the generalising solution is \textbf{seed-dominated}; the diameter filter only
        speeds up convergence (Reports~II--III).
  \item \textbf{There is a sharp, architectural distance wall at $\approx 3^L$ ($=9$ at $L=2$)},
        which moves with depth in a causal sweep and appears even on families seen $\sim 10^8$
        times in training --- steering the data cannot move it (Reports~III--IV).
  \item \textbf{Difficulty has two separable axes}: distance-vs-capacity and
        density-vs-training; the spectral gap is not a third axis on natural graphs (collinear with
        the diameter), though a purpose-built probe can isolate a bottleneck (Report~IV).
  \item \textbf{A similarity read-out doubles the reach to $2\cdot 3^L$} (a meet-in-the-middle of
        two $3^L$-hop neighbourhoods) but cannot cross a bottleneck (Report~IV).
  \item \textbf{The model computes distance-bounded matrix powering, not a bounded traversal}
        (DFS rejected; Report~V). The bridged-clique ``node budget'' ($\sim6$--$7$) that looked
        like a second capacity is a \textbf{soft, learned shortcut}: it moves with the read-out,
        \emph{dissolves once the structure is in training} (and the learned skill then transfers to
        an unseen bridge), and does not grow with depth.
\end{itemize}

\subsection{The new question: how does it learn a reasoning path, and from what data?}
\label{sec:question}

Report~V settled \emph{which} computation the model runs. The two facts that make this report
possible are the last one above: putting a structure in training can teach a \emph{within-capacity}
propagation skill the model otherwise skips, and that skill can \emph{transfer} to structures it
never saw. So the limiting factor for a given ``reasoning path'' (propagate a connection along a
particular shape) is often not the architecture but \emph{which data the model was trained on}.
This report asks that directly, through two lenses the transformer makes concrete:

\begin{itemize}\setlength\itemsep{0.15em}
  \item \textbf{Parallelism.} A transformer resolves every pair at once. Report~IV's
        confound-free probe already hinted that, beyond the capacity, \emph{several} parallel
        routes between two endpoints make a connection the model could not confirm over a single
        route. If real and robust, this is a clean statement about \emph{how} the model reasons:
        redundant evidence (more routes) is easier than a single thin thread, even at fixed
        distance. Report~IV measured this on one seed and on data partly saturated to $1.0$; here we
        redo it properly.
  \item \textbf{Iterated/repeated structure.} Report~V showed the model can be taught the
        single-bridge ``two cliques are one component iff the bridge is there'' decision. A natural
        next step is whether that local decision \emph{chains}: a sequence
        clique--bridge--clique--bridge--$\dots$ asks the model to find, at each stage, the node that
        carries the link to the next block --- a forced hand-off, repeated. This probes whether a
        learned local rule composes into a longer path of reasoning.
\end{itemize}

\subsection{The claims this report argues}
\label{sec:claims}

We keep the experiment set small and each one aimed at one claim. The intended spine:

\begin{enumerate}\setlength\itemsep{0.15em}
  \item \textbf{More parallel routes help} the model resolve a connection, \emph{even when it never
        saw such multi-route graphs in training} (parallelism aids the reasoning); and there is a
        measurable \textbf{sample cost} to learning it when the graphs \emph{are} shown.
  \item \textbf{Truncated routes in training hurt}: mixing in dead-end paths (a route that does not
        reach the far endpoint) confuses the connection skill rather than helping it.
  \item \textbf{Failing to handle the bridged cliques is ``never seen'', not ``too hard''}: like a
        too-dense ER graph it is an out-of-distribution gap, not an intrinsic difficulty --- shown
        by the fact that once the structure is in training the model handles it (Report~V), and we
        push this to repeated/varied versions.
  \item \textbf{A learned single bridge composes}: a model that learns one bridged clique can handle
        it \emph{repeated} (chains of bridged cliques) and with \emph{different block types} (e.g.\
        dense ER blocks instead of complete cliques), as long as the construction stays within the
        distance capacity ($\le 9$). \emph{(To be tested --- this is a hypothesis, not yet a
        result.)}
\end{enumerate}

\section{Setup and the data principle}
\label{sec:setup}

\paragraph{Model.} Unless noted we use the same base model as Reports~III--V: the RoBERTa-faithful
standard transformer with the linear read-out, $L=2$, single head, $d_{\mathrm{model}}=512$,
normalised-ReLU attention, trained online at $n=20$ (and $n=40$ where the canvas needs to be
larger) for $10^6$ steps. The input is the self-loop-augmented adjacency matrix $A$; the target is
the connectivity matrix $R$. Metrics are the usual exact-match, pairwise, and per-distance/per-pair
accuracies, plus, where relevant, the accuracy on a \emph{single designated pair} (the two
endpoints of a multi-route graph).

\paragraph{The data principle (new in this report).} Every experiment trains on a \emph{single,
explicitly named} graph distribution --- pure ER, cliques, bridged cliques, the multi-route
``multipath'' family, or a clearly stated and motivated combination of a few of them. We do
\emph{not} use the opaque uniform-over-nine-families ``mixed'' stream of Reports~III--V as a default
training set: a result is only worth reporting if the training distribution is simple enough to
state in one sentence and motivate. Where an earlier report's number came from the mixed stream and
we reuse it, we say so and treat it as a baseline, not as the clean training condition.

\paragraph{The two core families.}
\begin{itemize}\setlength\itemsep{0.15em}
  \item \textbf{Multipath} (the parallel family, from Report~IV): two terminals $a,b$ joined by $k$
        internally-disjoint paths each of length $\ell$, so the distance $\mathrm{dist}(a,b)=\ell$
        is fixed and only the number of routes $k$ varies (the effective resistance is $\ell/k$).
        The terminals are padded to a fixed degree and the rest of the canvas is filled sparsely, so
        $k$ is the only thing that changes. A \textbf{truncated} multipath replaces some routes with
        dead ends that do not reach $b$.
  \item \textbf{Bridged cliques} (the iterated family, from Report~V): two cliques (or two
        internally-dense blocks) of size $c$ joined by a single bridge edge --- one component. We
        extend it two ways here: a \textbf{chain} of cliques joined by successive bridges
        (clique--bridge--clique--$\dots$), and \textbf{thicker bridges} (more than one edge), keeping
        every cross distance $\le 9$ so the target is always within the matrix-power capacity.
\end{itemize}

\section{Plan of experiments}
\label{sec:plan}

The experiments fall into three threads, plus a short exploratory one.

\paragraph{Thread A --- Parallelism: the multipath family.} Redo Report~IV's parallel-route probe
properly (multi-seed, and only on the route/length cells that are actually informative), then ask
how many multipath graphs it takes to \emph{learn} the family, and whether truncated routes in
training hurt.

\paragraph{Thread B --- Asymmetric two-chains.} Explain a specific puzzle from Report~IV: a graph
split into two paths of very unequal length (e.g.\ $4$ and $36$) appeared \emph{easier} than a more
balanced split (e.g.\ $17$ and $23$). We analyse the existing data and run a controlled experiment,
reading it as a statement about how the model reasons along a path.

\paragraph{Thread C --- Iterated structure: chained and thickened bridged cliques.} Train on a clean
single-bridge family and test on repeated bridges (chains of cliques) and on thicker bridges and
different block types (dense ER blocks), to see whether a learned local link composes.

\paragraph{Thread D --- Trees (exploratory).} A short look at tree-structured graphs as another
reasoning-path shape.

\section{Results}
\label{sec:results}

\emph{(All subsections below are stubs: the runs are planned, not yet collected.)}

\subsection{Thread A.1 --- Parallel routes help, even unseen in training}
\label{sec:res-multipath-clean}
% SOURCE (eval-only): eval_multipath.py via scripts/r6_a1_eval.sbatch. Output:
%   runs/report6/multipath/n{N}_{dist}_seed{S}/multipath.json. Aggregate+plot locally with
%   plot_multipath.py -> runs/report6/report6_figs/multipath_{rescue,mechanism}_*.png.
% TRAININGS (clean single distributions; scripts/r6_a1_train.sbatch):
%   - ER: the prof's "never saw any path structure" framing. n20/n40 REUSE existing ER linear
%     checkpoints (repro_paper_n20_roberta 8 seeds; families_n40 4 seeds); n10/n64 trained fresh.
%   - path_union: sees SINGLE-path reach (sharp 3^L wall, k=1 fails beyond 9) but NEVER parallel
%     routes -> the clean rescue test. All sizes (n10/n20/n40/n64) trained fresh, 4 seeds.
%   - mixed: Report-IV reference only (NOT a clean condition; eval'd for comparison).
% SIZES n10/n20/n40/n64. FEASIBILITY need=2+k(ell-1)+2max(0,term_deg-k)<=n: the rescue (k>=2,
%   ell>9) needs room for >=2 long routes -> visible at n40 (k<=2-3) and CLEANLY at n64 (k<=4,
%   ell<=15); n10/n20 are within-capacity controls (k does not matter there). Report IV's clean
%   rescue (mixed n40) was: beyond cap k=1 ~0.55, k=2 ~0.93, k=3 ~1.00 at fixed distance.
% PRIMARY METRIC = the (a,b)=(s,t) PAIR accuracy (the graph IS the reasoning path; only the pair
%   matters). ALWAYS contextualise with the whole-matrix exact/pairwise (a correct pair with a
%   wrong matrix is itself a finding -- e.g. the pair is right because ONE route is fully resolved
%   while the others are wrong: it sees them connected via a SINGLE path). The per-route MECHANISM
%   (n_intact_hist: how many of the k routes are fully resolved) decides single-path vs genuine
%   multipath -- the heart of the analysis.
% CAPTION (rule 21): Model = base RoBERTa linear (L=2, 1 head, d512), {ER, path_union}-trained,
%   4 seeds (8 for n20 ER), n in {10,20,40,64}. Test = multipath, fixed distance ell, k routes,
%   >=400 graphs/cell. Metric = (s,t)-pair accuracy (+ whole-matrix exact), mean over seeds.
\paragraph{Question.} At a fixed endpoint distance, does adding parallel routes between $a$ and $b$
make the connection easier to resolve --- and does it, even for a model that never saw a multi-route
graph in training (pure ER)?

\emph{[To run --- data pending. Report~IV's single-seed version found that, beyond the capacity, one
route fails ($\approx0.55$) while two or three succeed; this redo establishes whether that holds
across seeds and identifies which route/length cells are informative rather than saturated.]}

\subsection{Thread A.2 --- The sample cost of learning the multipath connectivity}
\label{sec:res-multipath-samples}
% SOURCE: experiments2/train_multipath.py (trunc_frac=0) via scripts/r6_a2_samples.sbatch.
%   Output: runs/report6/multipath_train/n{N}_k{K}_ell{ELL}_clean_.../history.json with
%   val_pair_acc + val_exact vs step and steps_to{pair,exact}@{0.9,0.99,0.999} (samples=steps*1000).
% TRAINING: pure multipath stream (clean), one named distribution per run. Grid generalises over
%   n in {20,40,64}, routes k in {1..4}, length ell in {5,7,9,11,13}; 300k steps, eval every 2k,
%   4 seeds. Val = held-out clean multipath (same k,ell).
% PREDICTION (thesis 1, 2nd half): a measurable sample cost; does more routes (k) lower it, and
%   does longer ell (beyond capacity) raise it / cap the matrix? Read pair vs whole-matrix
%   separately -- the pair may be learned long before the matrix.
% CAPTION (rule 21): Model = base RoBERTa linear (L=2, 1 head, d512), multipath(k,ell)-trained,
%   4 seeds, n as labelled. Test = held-out clean multipath. Metric = steps (samples) to reach
%   the (s,t)-pair / whole-matrix exact threshold, mean over seeds.
\paragraph{Question.} When multipath graphs \emph{are} in training, how many does it take to learn
the connectivity exactly --- for the single $(a,b)$ pair, and for the whole matrix?

\emph{[To run --- data pending.]}

\subsection{Thread A.3 --- Truncated routes in training confuse the model}
\label{sec:res-multipath-truncated}
% SOURCE: experiments2/train_multipath.py (trunc_frac>0) via scripts/r6_a3_truncated.sbatch.
%   Output: runs/report6/multipath_train/n{N}_k{K}_ell{ELL}_trunc{f}_.../history.json (same dir
%   as A.2, so the clean trunc=0 baseline sits beside the truncated curves for plotting).
% TRAINING: a fraction f of training graphs are TRUNCATED -- n_full in {0..k-1} full routes + the
%   rest dead ends one hop short of t (labels correct; s-t connected iff n_full>=1). Eval on the
%   SAME clean multipath val as A.2. Configs n40 (k3,ell7) and (k2,ell9); f in {0.3,0.6}; baseline
%   f=0 is the matching A.2 run; 4 seeds.
% PREDICTION (thesis 2): truncated routes slow / degrade clean-multipath learning (the model must
%   verify a route completes, not just that one leaves a). Read pair AND matrix vs f.
% CAPTION (rule 21): Model = base RoBERTa linear (L=2, 1 head, d512), trained on multipath with a
%   fraction f of truncated routes, 4 seeds, n40. Test = clean held-out multipath (same k,ell).
%   Metric = (s,t)-pair / whole-matrix accuracy (and steps-to-threshold) vs f, mean over seeds.
\paragraph{Question.} If some training routes are dead ends (truncated), does that degrade what the
model learns about the genuine routes?

\emph{[To run --- data pending.]}

\subsection{Thread B --- Why an unbalanced two-chain split is easier}
\label{sec:res-asym-chains}
% ANALYSE the existing Report IV per-distance / per-seed two-chain data first (e.g. a 4+36 split
% vs a 17+23 split at n=40), then a controlled run. FRAMING: reasoning along a path -- a long path
% next to a very short one may be effectively a single long reach, whereas two medium paths put two
% capacity-boundary reaches in the same graph. Read from the data, do not assume the mechanism.
% CAPTION (rule 21): Model = base RoBERTa linear (L=2, 1 head, d512), [training set as run], N seeds.
%   Test = two-chain graphs at the chosen length splits, #graphs. Metric = exact-match / per-pair
%   accuracy by split, mean over seeds.
\paragraph{Question.} Report~IV suggested that splitting $n$ nodes into two paths of very unequal
length is \emph{easier} than a balanced split. Why, and what does it say about how the model reasons
along a path?

\emph{[To analyse the existing data, then run --- pending.]}

\subsection{Thread C.1 --- Chained bridged cliques: does a learned bridge compose?}
\label{sec:res-chain-cliques}
% TRAINING: a single bridged-clique family (clean, as in Report V sec 5.6, bridged in stream).
% TEST: chains clique-bridge-clique-bridge-... (repeated forced hand-offs), within distance <= 9.
% PREDICTION (thesis 4): if the single bridge is learned, the chained version is handled too (the
% local "iff + find the bridge node" rule composes). FAILURE would mean the skill is shape-specific.
% CAPTION (rule 21): Model = base RoBERTa linear, trained on single bridged cliques, N seeds.
%   Test = chained bridged cliques (#cliques in the chain), #graphs. Metric = cross-block / exact-match
%   by number of links in the chain, mean over seeds.
\paragraph{Question.} A single bridged clique is a learnable ``iff'' decision (Report~V). Does it
\emph{chain} --- can the model carry the link across clique--bridge--clique--$\dots$, finding at each
stage the node that hands off to the next block?

\emph{[To run --- data pending.]}

\subsection{Thread C.2 --- Thicker bridges and different blocks}
\label{sec:res-thick-bridges}
% VARIANTS, all kept within distance <= 9: (i) bridges with MORE THAN ONE edge (does redundancy in
% the bridge make it easier? ties to Thread A parallelism); (ii) BLOCKS that are dense ER blocks
% (internal edge prob ~0.6) joined by a bridge, instead of complete cliques (the held-out "bridged
% dense blocks" of Report V sec 5.6) -- easier or harder than cliques? PREDICTION (thesis 3-4):
% within capacity these are learnable; thicker bridges may be easier (more routes), ER blocks
% comparable once the block density is in-distribution.
% CAPTION (rule 21): Model = base RoBERTa linear, [training set as run], N seeds.
%   Test = bridged cliques with bridge thickness t / dense ER blocks, sweep size c, #graphs.
%   Metric = cross-block pair accuracy vs size, by bridge thickness / block type, mean over seeds.
\paragraph{Question.} Does a thicker bridge (more than one edge) make the hand-off easier, and do
dense ER blocks behave like complete cliques --- all kept within the distance capacity?

\emph{[To run --- data pending.]}

\subsection{Thread D --- Trees (exploratory)}
\label{sec:res-trees}
% A short exploration of tree-structured graphs as another reasoning-path shape. Scope TBD;
% likely a single clean tree family. Keep brief unless it yields a clear result.
\paragraph{Question.} How does the model handle tree-structured connectivity, as a third
reasoning-path shape beside parallel routes and chained bridges?

\emph{[Exploratory --- scope and runs pending.]}

\section{Verdict (to write last)}
\label{sec:verdict}
% Synthesise around the claims of sec 1.3: parallelism helps (and unseen-in-training);
% truncated routes hurt; bridged-clique failure is "unseen" not "hard"; a learned bridge composes.
% Tie back to Report V (distance-bounded matrix powering; the node budget is a removable data prior).
% State clearly, for each claim, whether the evidence supports it. Keep honesty: behavioural
% evidence, the seed lottery, the ER/density confounds.
\emph{[To write last, once Threads A--C are in.]}

\end{document}
